\documentclass[11pt]{article}

\usepackage[T1]{fontenc}
\usepackage[margin=0.8in]{geometry}
\usepackage{booktabs}
\usepackage{enumitem}
\usepackage{hyperref}
\usepackage{longtable}
\usepackage{xcolor}

\hypersetup{colorlinks=true,linkcolor=blue,urlcolor=blue}
\setlist{nosep,leftmargin=*}
\newcommand{\good}{\textcolor{green!45!black}{PASS}}
\newcommand{\bad}{\textcolor{red!70!black}{FAIL}}
\newcommand{\pending}{\textcolor{orange!80!black}{PENDING}}

\title{Unified Dynamic-Manipulation Generator Repair\\Technical Handoff}
\author{Codex repair session for \texttt{zl664/dataset\_generation}}
\date{18 July 2026}

\begin{document}
\maketitle

\section{Purpose of this handoff}

This document records the exact state of the unified robotics-dataset repair so
that work can continue safely in another chat.  It distinguishes committed
code, uncommitted work, sealed evidence, partial outputs, automated passes, and
qualitative failures.  The most important current conclusion is:

\begin{quote}
The P0b/P0c camera and reviewability repair is successful, but the current F1
controller is not acceptable as a dynamic reaching catch.  It initializes the
robot at the interception pose and changes only the gripper command.  The F1
preview must therefore remain blocked even though automated free-contact QC
passes.
\end{quote}

No hundreds-of-hours generation should be launched from the current state.

\section{Repository, branch, and exact Git state}

\begin{itemize}
  \item Repository: \path{/gpfs/radev/project/sous/zl664/dataset_generation}
  \item Branch: \texttt{codex/unified-generator-repair}
  \item Current committed HEAD: \texttt{ac56d64} (\texttt{Repair passive review camera evidence})
  \item Baseline/checkpoint commits, newest first:
    \begin{itemize}
      \item \texttt{ac56d64} --- P0b/P0c review-camera repair and visibility-QC v3.
      \item \texttt{d1a56f3} --- fixed-seed physics and review blockers.
      \item \texttt{0a6495b} --- randomized-scene physics and review evidence.
      \item \texttt{e613fc8} --- persisted wall-rebound review evidence.
      \item \texttt{5ed2581} --- unified review-only generation pipeline.
      \item \texttt{fc65503} --- embodiment-correction baseline.
    \end{itemize}
\end{itemize}

At handoff time the worktree is intentionally dirty.  The following files are
an \emph{uncommitted and not yet fully tested} memory-repair prototype:

\begin{itemize}
  \item modified: \path{src/dynamic_robot_dataset/common/source_preview.py}
  \item modified: \path{tests/unit/test_source_preview.py}
  \item new: \path{src/dynamic_robot_dataset/common/source_preview_prepare_worker.py}
\end{itemize}

The last focused pytest command for these files was interrupted by the user
after 3.9 seconds.  It produced no trustworthy pass/fail result, and no pytest
or generator process was left running.  The next chat must inspect and test
these changes before committing them.

\section{What the collaborator's ``harmonized'' directory contains}

The collaborator clarified that the operative framework is:

\path{/gpfs/radev/project/sous/zss8/dataset-generation/harmonized/}

Important files are:

\begin{itemize}
  \item \texttt{harmonize.py}: entry point for plan, scaffold, convert, generate,
        validate, and status.
  \item \texttt{taxonomy.py}: collaborator-side family/subfamily registry.
  \item \texttt{converters/}: conversion adapters into the canonical format.
\end{itemize}

The collaborator reported that only the v3-ish converter is currently wired
and conversion must run in the framework virtual environment.  Their
\texttt{generate} command does not supply a different physical F1 controller;
it delegates to this repository's \texttt{source\_mujoco} backend.  Therefore
the static-arm F1 problem must be repaired here, not in
\texttt{harmonize.py}.

The directory
\path{/gpfs/radev/scratch/sous/zss8/dynamic_manip_corpus/}
primarily exposed JSON/configuration artifacts during inspection; it was not
the source of the generator implementation.  All external generator and
RoboCasa trees were treated as read-only dependencies.

\section{Unified architecture already implemented}

The branch now contains a single orchestration/schema/QC path rather than a
separate conversion-only runtime.  Implemented elements include:

\begin{itemize}
  \item a 20-leaf corpus registry and backend-capability registry;
  \item a serializable, hash-bound \texttt{SourceScenarioSpec};
  \item immutable \texttt{plan-run}, deterministic \texttt{run-shard}, sealed
        \texttt{finalize-run}, bounded \texttt{generate}, and fixed
        \texttt{review-suite} paths;
  \item canonical Arrow schemas, two synchronized 832-by-480 MP4 streams at
        30 Hz, frame timestamps, simulation timestamps, and synchronization
        errors;
  \item actual eight-dimensional actuator actions for Panda/Robotiq rather
        than desired 15-joint diagnostics;
  \item mutation-boundary enforcement: runtime callbacks may change
        \texttt{data.ctrl}, not robot/object state, latches, welds, mocap,
        applied forces, or mutable model physics;
  \item strict contact, penetration, energy, rebound, grasp, visibility,
        RoboCasa-clearance, source-hash, and saved-artifact replay checks;
  \item immutable event strips and a hash-bound pending human-review ledger.
\end{itemize}

``Unified'' here means one registry, schema, writer, evaluator/QC interface,
and orchestration layer across validated backends.  It does not claim that
MuJoCo, Genesis fluid, and deformable solvers share identical internal
physics.

\section{Current leaf activation status}

The fixed review plan contains exactly 20 leaves times 6 cases = 120 logical
episodes and 240 requested canonical MP4s.  The immutable plan has 48
executable cases and assigns zero quota to 12 unfinished leaves.

\begin{longtable}{lll}
\toprule
Leaves & Current execution state & Notes \\
\midrule
\endhead
P0a--P0d & review-executable & passive rigid physics, no robot by design \\
F1a--F1d & review-executable but qualitatively blocked & free-contact catch path exists; reaching motion is defective \\
F2a--F2f & blocked & direct/rebound/ramp/multisurface adapters unfinished \\
F3a, F3b, F3d & blocked & platform/rolling/moving-base adapters unfinished \\
F3c & blocked & Genesis current-and-pickup scene not implemented \\
D1, D2 & blocked & frictional cloth/rope grasp not implemented \\
\bottomrule
\end{longtable}

The permanently synthetic \texttt{native\_mujoco} backend remains blocked.
Legacy state-rewritten or latch-assisted catches remain in a separate
\texttt{legacy\_assisted} quarantine tier and must not count toward training
or generated hours.

\section{Committed v8 camera and reviewability repair}

Commit \texttt{ac56d64} introduced:

\begin{itemize}
  \item runtime profile \texttt{source-mujoco-rigid-review-2026-07-v8};
  \item backend version \texttt{0.8.0-review};
  \item visibility-QC schema \texttt{source-mujoco-visibility-qc/v3};
  \item unchanged compiled-scenario schema \texttt{/v3};
  \item preserved external source and RoboCasa hashes.
\end{itemize}

The exact P0b main-camera repair is:

\begin{verbatim}
position = (1.20, -1.15, support_height + 1.15)
target   = (0.05, -0.10, support_height + 0.85)
fovy     = 58 degrees
\end{verbatim}

The exact wall-only P0c secondary-camera repair is:

\begin{verbatim}
position = (-0.05, -1.70, support_height + 0.76)
target   = (-0.05,  0.00, support_height + 0.76)
fovy     = 58 degrees
\end{verbatim}

P0c table-rebound secondary-camera serialization was deliberately left
unchanged.

Visibility-QC v3 now requires all initial/apex/key-event/final checkpoints to
be visible in P0b's main stream.  For P0c wall rebounds it requires those four
checkpoints in the secondary stream and independently recomputes
\(\lvert \hat v_{camera}\cdot\hat n_{wall}\rvert\leq 0.05\).  Malformed task,
camera, fixture, half-size, or quaternion evidence fails closed.

\section{Test results at committed v8}

After commit \texttt{ac56d64}, the complete test matrix passed:

\begin{itemize}
  \item 406 non-integration tests passed; 51 integration tests deselected.
  \item All 51 integration tests passed in graphics-safe isolated batches:
        22 + 4 + 7 + 5 + 6 + 4 + 3.
  \item \texttt{compileall} passed.
  \item \texttt{git diff --check} passed.
\end{itemize}

Graphics tests were intentionally batched because a monolithic EGL lifecycle
has previously aborted after many rendered tests.  The isolated results are
the trustworthy baseline.

\section{Immutable review-plan identity}

The fresh v8 plan is:

\path{/gpfs/radev/project/sous/zl664/dataset_generation_runs/review_suite_20260718_v8_ac56d64}

Its identities are unchanged from v7, proving that the same difficult fixed
cases and seeds were preserved:

\begin{itemize}
  \item plan SHA-256:
        \texttt{c11a8065cd87099ed4379198d936e5bc9e9c82236294d5d8a4fd1399080422a9}
  \item request-ledger SHA-256:
        \texttt{c5b87826a108cb1a869c8373a2ef21192081f4a8bead4dfbb65a1c91ea09b1fe}
  \item suite ID: \texttt{review-suite-197cfe41997cb8e5abc04b24}
\end{itemize}

\section{Generated datasets and whether they are usable}

\subsection{Sealed v7 reference evidence}

\path{/gpfs/radev/project/sous/zl664/dataset_generation_runs/review_active8_fixed6_v7_d1a56f3}

This contains 48 episodes, 96 MP4s, 96 event-strip PNGs, exact finalization,
and zero automated hard failures under the then-current contract.  It is now
historical failed review evidence because P0b-review-02 was clipped in main
and P0c wall recoil was visually ambiguous.  Visibility-QC v3 intentionally
does not accept its v2 evidence.

\subsection{Sealed v8 P0b/P0c repair bundle --- usable review evidence}

\path{/gpfs/radev/project/sous/zl664/dataset_generation_runs/review_p0bc_fixed6_v8_ac56d64}

Status: \good.  It contains 12 episodes and 24 MP4s.  All 12 cases passed
strict saved-artifact replay.  Independent qualitative inspection also passed
all cases.  In particular:

\begin{itemize}
  \item P0b-review-02 now shows the complete rise, apex, and descent in both
        streams with comfortable margins.
  \item P0c-review-01/03/05 secondary streams clearly show approach, wall
        contact, and wall-normal separation.
  \item no visible tunneling, sticking, unexplained jumps, floating fixtures,
        task-volume background intersections, or material occlusion were found.
\end{itemize}

These are passive-physics leaves.  Their declarations use
\texttt{embodiment=no\_robot}; the absence of a Franka arm is intentional.

\subsection{Sealed one-case F1a Franka preview --- automated pass, qualitative fail}

\path{/gpfs/radev/project/sous/zl664/dataset_generation_runs/preview_f1a_franka_v8_ac56d64}

Direct media paths:

\begin{itemize}
  \item main MP4:
    \path{/gpfs/radev/project/sous/zl664/dataset_generation_runs/preview_f1a_franka_v8_ac56d64/videos/observation.images.main/chunk-000/file-000000.mp4}
  \item secondary MP4:
    \path{/gpfs/radev/project/sous/zl664/dataset_generation_runs/preview_f1a_franka_v8_ac56d64/videos/observation.images.secondary/chunk-000/file-000000.mp4}
  \item main event strip:
    \path{/gpfs/radev/project/sous/zl664/dataset_generation_runs/preview_f1a_franka_v8_ac56d64/reviews/event_strips/F1a/F1a-review-00/main.png}
\end{itemize}

Automated physics, free-contact, retention, media, and strict QC passed.  The
user correctly identified a decisive qualitative failure: the arm waits at
the catch point instead of reaching for the ball.  This output must not be
used as accepted F1 review evidence.

\subsection{Interrupted/unsealed outputs --- do not use}

The following runs were killed by the Slurm memory cgroup and remain partial,
unsealed interruption evidence:

\begin{itemize}
  \item full active-eight attempt, 3 committed episodes / 6 MP4s:
    \path{/gpfs/radev/project/sous/zl664/dataset_generation_runs/review_active8_fixed6_v8_ac56d64}
  \item F1a six-case attempt, 3 committed episodes / approximately 6 complete
        MP4s plus a possible partially written next file:
    \path{/gpfs/radev/project/sous/zl664/dataset_generation_runs/review_F1a_fixed6_v8_ac56d64}
\end{itemize}

Neither contains a final seal or physics-failure receipt.  They are
interruptions, not outcome mismatches, but should not be published or treated
as review-complete.

\section{A40 GPU and the actual performance blocker}

The current environment exposes an NVIDIA A40 with 46,068 MiB.  During an
active render worker, approximately 528 MiB was allocated.  The GPU was not
the cause of the slow/failed full run.

The actual blocker is eager native-memory retention in the parent process:

\begin{itemize}
  \item \texttt{execute\_source\_review\_preview} prepared every declaration
        in the parent before launching isolated render shards.
  \item Each F1 declaration compiles provisional and final MuJoCo/RoboCasa
        models.  Native allocations remain resident after Python objects are
        deleted.
  \item Same-process preparation of 12 F1a/F1b cases grew RSS from 381 MB to
        2.04 GB while Python heap stayed near 22 MB.
  \item Six F1a cases still grew the parent heavily even with
        \texttt{malloc\_trim}.
  \item The full parent was killed at about 2.58 GB anonymous RSS while a
        render child was active.
  \item The F1a-only parent was later killed at about 1.61 GB anonymous RSS
        while its child was active.
  \item In contrast, preparing six F1a declarations one-per-subprocess kept
        parent RSS flat at 40.5 MB; returned JSON was only 12--17 KB per case.
\end{itemize}

The compile chain is:

\begin{itemize}
  \item \texttt{source\_execution.py}, declaration preparation;
  \item \texttt{source\_spec.py}, source-model compilation;
  \item \texttt{model.py}, provisional R1 model;
  \item \texttt{model.py}, final model.
\end{itemize}

The uncommitted prototype fixes this by preparing exactly one declaration per
fresh owned subprocess, returning only canonical JSON, validating case/UUID/
index/hash/commit identity in the parent, and letting process exit reclaim all
native memory.  Existing one-shard-per-process execution remains unchanged.

\section{Static Franka catch: diagnosis and required repair}

The user's observation was verified quantitatively from the sealed F1a
preview Parquet:

\begin{itemize}
  \item all seven arm actuator-command peak-to-peak ranges are exactly zero;
  \item only the Panda finger command changes, from 255 to approximately 131.5;
  \item maximum observed arm-joint displacement is only 0.0074 rad and is
        passive contact/servo deflection, not commanded reaching;
  \item \texttt{source\_mujoco/backend.py} sets both
        \texttt{initial\_arm\_command} and
        \texttt{intercept\_arm\_command} to the same IK solution and initializes
        robot qpos at that intercept.
\end{itemize}

This behavior was introduced in commit \texttt{5ed2581}; it was not inherited
from the earlier generator.

The original Franka controller at
\path{legacy_sources/zss8/dataset-generation/demo_mujoco_arm_gripper/franka_catch/controller.py}
does show reaching.  Its transferable planning intent is:

\begin{itemize}
  \item solve an intercept pose;
  \item solve a ready waypoint 45 mm above it, seeded from the intercept IK;
  \item minimum-jerk ready-to-intercept motion;
  \item reach end at ballistic event minus 0.055 s;
  \item duration at least 0.18 s and long enough for joint limits;
  \item for the current 0.44 s event, approximately 0.205--0.385 s;
  \item an optional visible post-catch lift/retract of roughly
        \((-0.15,0,+0.26)\) m over at least 0.55 s.
\end{itemize}

The original execution itself must \emph{not} be copied: it directly rewrites
robot qpos/qvel each step and rewrites/holds the ball after capture.  The
collaborator's converter correctly labels those episodes
\texttt{assisted\_contact}.  The safe adaptation must:

\begin{enumerate}
  \item solve ready and intercept IK only during initialization;
  \item initialize once at the ready waypoint;
  \item move ready-to-intercept using minimum-jerk commands applied only via
        \texttt{data.ctrl};
  \item use a bounded closure rather than the original instantaneous close;
  \item never rewrite or latch the ball;
  \item preserve measured success/failure and the same fixed review seeds;
  \item serialize the actual arm and closure phases honestly in
        \texttt{SourceScenarioSpec}.
\end{enumerate}

This is a profile-changing repair and should become v9 (and likely backend
\texttt{0.9.0-review}), not a silent edit to v8.

\section{Exact recommended next steps for the new chat}

\begin{enumerate}
  \item Read this file, run \texttt{git status}, and preserve the three
        uncommitted memory-fix files.  Do not overwrite or delete the partial
        run directories.
  \item Finish and test the isolated declaration-preparation worker:
    \begin{itemize}
      \item run \texttt{tests/unit/test\_source\_preview.py};
      \item test nonzero child exit, malformed/non-mapping JSON, and identity
            mismatch;
      \item run a real six-case F1a preparation RSS probe;
      \item run the full non-integration suite and relevant integration tests;
      \item commit the memory repair separately.
    \end{itemize}
  \item Implement the v9 ctrl-only reaching controller:
    \begin{itemize}
      \item ready IK target = intercept + \((0,0,0.045)\) m;
      \item initialize at ready, not intercept;
      \item compute feasible minimum-jerk timing and reject infeasible IK or
            joint-limit violations;
      \item coordinate bounded closure with arrival and retain real contact;
      \item update \texttt{source\_spec.py} so actuator phases show the reach;
      \item add QC requiring nonzero arm command travel and arrival evidence;
      \item add 600/1200 Hz same-seed agreement and mutation-boundary tests.
    \end{itemize}
  \item Generate only \texttt{F1a-review-00} first.  Qualitatively verify that
        the arm visibly reaches, the ball remains free-contact, bilateral
        support occurs, and the final grasp is retained.
  \item Rerun the same six F1a seeds, then F1b--F1d, without replacing difficult
        seeds.  Preserve actual failures and negative branches.
  \item After the memory fix is committed, create a fresh full output whose
        directory includes the new commit/profile.  Do not resume an ac56d64
        plan after changing generator Git identity.
  \item Run standalone \texttt{validate --strict-all}, external strict QC into
        a separate report directory, and qualitative two-view/event-strip
        review.
  \item Leave the human-review ledger pending; do not fabricate approvals.
  \item Only after all six reviews for a leaf pass should it enter a
        100-episode pilot and two-shard interruption/resume rehearsal, followed
        by the 10-hour and 100-hour gates.
\end{enumerate}

\section{Scale-readiness warning}

The system is not ready for hundreds of hours.  Remaining hard gates include:

\begin{itemize}
  \item dynamic reaching repair for F1;
  \item formal human review;
  \item 12 blocked leaves with zero quota;
  \item 100-episode pilot with zero hard failures;
  \item two-shard interruption/resume rehearsal;
  \item 10-unique-hour and 100-unique-hour gates.
\end{itemize}

The fixed six-case review matrix intentionally correlates scene order and
outcome branch.  A prior statistics report measured background/outcome
Cramer's \(V\approx0.53\).  This is acceptable for deterministic acceptance
testing but must not be copied as the scale-generation randomization policy.

\section{Bottom line}

The passive P0b/P0c repair is good and reproducibly validated.  The unified
writer, schema, orchestration, physics/QC, RoboCasa binding, and fixed-review
infrastructure are substantially implemented.  Two issues now dominate:

\begin{enumerate}
  \item declaration preparation must be isolated to prevent native-memory OOM;
  \item F1 must perform a genuine actuator-only reaching catch rather than a
        stationary pre-positioned interception.
\end{enumerate}

Resolve those two items before generating more F1 review media or considering
large-scale data production.

\end{document}
